\documentclass[11pt,a4paper]{article}
\usepackage[utf8]{vietnam}
\usepackage{amsmath,amssymb}
\usepackage{tabularx}
\usepackage{booktabs}
\usepackage{multirow}
\usepackage[top=1.0cm, bottom=1.0cm, left=1.2cm, right=1.2cm]{geometry}
\usepackage{fancyhdr}

\pagestyle{fancy}
\fancyhf{}
\rfoot{\small\textit{Trang \thepage}}
\lfoot{\small\textit{Hướng dẫn chấm -- Đề 2 (Toán 7, Giữa HK I)}}
\renewcommand{\headrulewidth}{0pt}

\newcolumntype{Y}{>{\centering\arraybackslash}X}

\begin{document}

\begin{center}
  \textbf{\large ĐÁP ÁN VÀ HƯỚNG DẪN CHẤM -- ĐỀ 2}\\[1mm]
  \textbf{\large KIỂM TRA GIỮA HỌC KÌ I -- MÔN TOÁN 7 -- NĂM HỌC 2026 -- 2027}
\end{center}

\vspace{1mm}
\noindent\textbf{\large PHẦN I. (3,0 điểm) Mỗi câu trả lời đúng được 0,25 điểm.}
\vspace{1mm}

\noindent
\begin{tabularx}{\linewidth}{|c|*{12}{Y|}}
\hline
\textbf{Câu} & 1 & 2 & 3 & 4 & 5 & 6 & 7 & 8 & 9 & 10 & 11 & 12 \\ \hline
\textbf{Đáp án} & \textbf{C} & \textbf{B} & \textbf{A} & \textbf{D} & \textbf{C} & \textbf{B} & \textbf{D} & \textbf{A} & \textbf{C} & \textbf{B} & \textbf{D} & \textbf{A} \\ \hline
\end{tabularx}

\vspace{2.5mm}
\noindent\textbf{\large PHẦN II. (2,0 điểm)}
\vspace{0.5mm}

\noindent\textit{Quy tắc chấm: Đúng 1 ý được 0,1đ; đúng 2 ý được 0,25đ; đúng 3 ý được 0,5đ; đúng 4 ý được 1,0đ.}

\vspace{1mm}
\noindent
\begin{tabularx}{\linewidth}{|c|c|c|X|}
\hline
\textbf{Câu} & \textbf{Ý} & \textbf{Đáp án} & \textbf{Giải thích tóm tắt} \\ \hline
\multirow{4}{*}{\textbf{13}} & a) & \textbf{Đúng} & 7 không là bình phương của số hữu tỉ nên $\sqrt{7}$ và $x$ là số vô tỉ. \\ \cline{2-4}
& b) & \textbf{Sai} & $\dfrac{7}{11} = 0,(63)$, không phải $0,6(3)$. \\ \cline{2-4}
& c) & \textbf{Sai} & $x \approx -2,646 < z = -2,6 \Rightarrow x < z < t$. Khẳng định $z < x < t$ là sai. \\ \cline{2-4}
& d) & \textbf{Đúng} & $m + 1 = 2,5 \Rightarrow m = 1,5$ hoặc $m + 1 = -2,5 \Rightarrow m = -3,5$; tổng $= -2$. \\ \hline
\multirow{4}{*}{\textbf{14}} & a) & \textbf{Đúng} & $\widehat{M}_1$ và $\widehat{M}_2$ là hai góc kề bù (chung cạnh, hai cạnh kia đối nhau). \\ \cline{2-4}
& b) & \textbf{Sai} & $\widehat{M}_1$ và $\widehat{N}_1$ là cặp góc đồng vị, không phải so le trong. \\ \cline{2-4}
& c) & \textbf{Đúng} & $\widehat{M}_1 = 180^\circ - 115^\circ = 65^\circ = \widehat{N}_1$, hai góc đồng vị bằng nhau nên $a \parallel b$. \\ \cline{2-4}
& d) & \textbf{Đúng} & $\widehat{M}_3 = \widehat{M}_1 = 65^\circ$; tia phân giác tạo góc $32,5^\circ \Rightarrow \widehat{MPN} = 32,5^\circ$ (so le trong). \\ \hline
\end{tabularx}

\vspace{2.5mm}
\noindent\textbf{\large PHẦN III. (2,0 điểm) Mỗi câu đúng được 0,5 điểm.}
\vspace{1mm}

\noindent
\begin{tabularx}{\linewidth}{|c|c|X|}
\hline
\textbf{Câu} & \textbf{Đáp án} & \textbf{Hướng dẫn giải tóm tắt} \\ \hline
\textbf{15} & \textbf{1} & $3\sqrt{16} - \sqrt{121} = 3 \cdot 4 - 11 = 1$. \\ \hline
\textbf{16} & \textbf{2} & $\dfrac{2}{3}x = \dfrac{1}{2} + \dfrac{5}{6} = \dfrac{4}{3} \Rightarrow x = 2$. \\ \hline
\textbf{17} & \textbf{55} & $a \parallel b \Rightarrow \widehat{K}_2 = \widehat{H}_2 = 125^\circ$ (đồng vị); $\widehat{K}_1 = 180^\circ - 125^\circ = 55^\circ$ (kề bù). \\ \hline
\textbf{18} & \textbf{35} & Cạnh tấm bìa: $\sqrt{75} \approx 8,66\text{ cm}$. Ruy băng: $4\sqrt{75} \approx 34,64 \approx 35\text{ cm}$. \\ \hline
\end{tabularx}

\vspace{2.5mm}
\noindent\textbf{\large PHẦN IV. Tự luận (3,0 điểm)}
\vspace{1mm}

\noindent
\begin{tabularx}{\linewidth}{|c|c|X|c|}
\hline
\textbf{Câu} & \textbf{Ý} & \textbf{Nội dung lời giải} & \textbf{Điểm} \\ \hline
\multirow{4}{*}{\textbf{19}} & \multirow{2}{*}{a)} & $\dfrac{4}{9} \cdot \dfrac{13}{5} - \dfrac{4}{9} \cdot \dfrac{8}{5} + \dfrac{5}{9} = \dfrac{4}{9} \cdot \left(\dfrac{13}{5} - \dfrac{8}{5}\right) + \dfrac{5}{9}$ & 0,25 đ \\ \cline{3-4}
& & $= \dfrac{4}{9} \cdot 1 + \dfrac{5}{9} = \dfrac{9}{9} = 1$. & 0,25 đ \\ \cline{2-4}
& \multirow{2}{*}{b)} & $\left(2x + \dfrac{1}{3}\right)^2 = \dfrac{25}{9} \Rightarrow 2x + \dfrac{1}{3} = \dfrac{5}{3}$ hoặc $2x + \dfrac{1}{3} = -\dfrac{5}{3}$. & 0,25 đ \\ \cline{3-4}
& & Với $2x + \dfrac{1}{3} = \dfrac{5}{3} \Rightarrow x = \dfrac{2}{3}$; với $2x + \dfrac{1}{3} = -\dfrac{5}{3} \Rightarrow x = -1$. Vậy $x \in \left\{\dfrac{2}{3}; -1\right\}$. & 0,25 đ \\ \hline

\multirow{4}{*}{\textbf{20}} & \multirow{2}{*}{a)} & Lớp 7A mượn: $240 \cdot \dfrac{1}{3} = 80\text{ (quyển)}$. Còn lại: $240 - 80 = 160\text{ (quyển)}$.\newline Lớp 7B mượn: $160 \cdot \dfrac{3}{8} = 60\text{ (quyển)}$. & 0,25 đ \\ \cline{3-4}
& & Sau hai lớp mượn, thư viện còn lại: $160 - 60 = 100\text{ (quyển)}$. & 0,25 đ \\ \cline{2-4}
& \multirow{2}{*}{b)} & Số truyện mua thêm: $100 \cdot \dfrac{2}{5} = 40\text{ (quyển)}$.\newline Giá mỗi quyển sau giảm 12\%: $50\,000 \cdot (100\% - 12\%) = 44\,000\text{ (đồng/quyển)}$. & 0,25 đ \\ \cline{3-4}
& & Số tiền nhà trường phải trả: $40 \cdot 44\,000 = 1\,760\,000\text{ (đồng)}$. & 0,25 đ \\ \hline

\multirow{4}{*}{\textbf{21}} & \multirow{2}{*}{a)} & Ta có $\widehat{ABn'} = 180^\circ - \widehat{ABn} = 180^\circ - 70^\circ = 110^\circ$ (hai góc kề bù). & 0,25 đ \\ \cline{3-4}
& & Mà $\widehat{mAB} = \widehat{ABn'} = 110^\circ$, hai góc này ở vị trí so le trong nên $mm' \parallel nn'$. & 0,25 đ \\ \cline{2-4}
& \multirow{2}{*}{b)} & Vì $d \perp mm'$ mà $mm' \parallel nn'$ nên $d \perp nn'$ (từ vuông góc đến song song). & 0,25 đ \\ \cline{3-4}
& & Ta có $\widehat{mAC} = 90^\circ$ (do $d \perp mm'$). Vì tia $AC$ nằm giữa hai tia $Am, AB$ nên:\newline $\widehat{BAC} = \widehat{mAB} - \widehat{mAC} = 110^\circ - 90^\circ = 20^\circ$. & 0,25 đ \\ \hline
\end{tabularx}

\end{document}
